\documentclass{ximera}

\input{../../../preamble.tex}


\definecolor{fvdnavy}{HTML}{071D43}
\definecolor{fvdcyan}{HTML}{20BFC6}
\definecolor{fvdcyanlight}{HTML}{EFFCFB}
\definecolor{fvdmagenta}{HTML}{F10D69}
\definecolor{fvdmagentalight}{HTML}{FFF1F7}
\definecolor{fvdtext}{HTML}{163044}





%=================================================
% Pedagogische omgevingen
% PDF: tcolorbox
% HTML: learning-box
%=================================================

\newenvironment{voorkennis}
{%
    \ifdefined\HCode
        \HCode{%
<section class="learning-box learning-box--prior">
<div class="learning-box__header">
<span class="learning-box__icon" aria-hidden="true"></span>
<span class="learning-box__title">Voorkennis</span>
</div>
<div class="learning-box__content">
        }%
        \begin{itemize}
    \else
       \begin{tcolorbox}[
    enhanced,
    breakable,
    colback=fvdcyanlight,
    colframe=fvdcyan,
    coltext=fvdtext,
    boxrule=0.5pt,
    leftrule=4pt,
    arc=4mm,
    outer arc=4mm,
    left=7mm,
    right=6mm,
    top=5mm,
    bottom=5mm,
    before skip=6mm,
    after skip=6mm,
    title=Voorkennis,
    coltitle=fvdcyan,
    fonttitle=\bfseries\Large,
    attach title to upper={\par\smallskip},
    boxed title style={
        empty,
        left=0mm,
        right=0mm,
        top=0mm,
        bottom=0mm
    },
    overlay={
        \draw[
            fvdcyan,
            line width=0.5pt
        ]
        ([xshift=42mm,yshift=-13mm]frame.north west)
        --
        ([xshift=-7mm,yshift=-13mm]frame.north east);
    }
]
        \begin{itemize}
    \fi
}
{%
    \end{itemize}
    \ifdefined\HCode
        \HCode{%
</div>
</section>
        }%
    \else
        \end{tcolorbox}
    \fi
}


\newenvironment{leerdoelen}
{%
    \ifdefined\HCode
        \HCode{%
<section class="learning-box learning-box--goals">
<div class="learning-box__header">
<span class="learning-box__icon" aria-hidden="true"></span>
<span class="learning-box__title">Leerdoelen</span>
</div>
<div class="learning-box__content">
        }%
        \begin{itemize}
    \else
\begin{tcolorbox}[
    enhanced,
    breakable,
    colback=fvdmagentalight,
    colframe=fvdmagenta,
    coltext=fvdtext,
    boxrule=0.5pt,
    leftrule=4pt,
    arc=4mm,
    outer arc=4mm,
    left=7mm,
    right=6mm,
    top=5mm,
    bottom=5mm,
    before skip=6mm,
    after skip=6mm,
    title=Leerdoelen,
    coltitle=fvdmagenta,
    fonttitle=\bfseries\Large,
    attach title to upper={\par\smallskip},
    boxed title style={
        empty,
        left=0mm,
        right=0mm,
        top=0mm,
        bottom=0mm
    },
    overlay={
        \draw[
            fvdmagenta,
            line width=0.5pt
        ]
        ([xshift=42mm,yshift=-13mm]frame.north west)
        --
        ([xshift=-7mm,yshift=-13mm]frame.north east);
    }
]
        \begin{itemize}
    \fi
}
{%
    \end{itemize}
    \ifdefined\HCode
        \HCode{%
</div>
</section>
        }%
    \else
        \end{tcolorbox}
    \fi
}


\addPrintStyle{../../..}

\title{Van data naar puntenwolk}
\author{Ingmar Herreman}

\begin{abstract}
Wanneer twee grootheden samen worden gemeten, ontstaat een nieuwe vraag:
\emph{hangt de ene grootheid samen met de andere?}

Een tabel met meetwaarden bevat alle afzonderlijke waarnemingen,
maar maakt een patroon niet altijd zichtbaar. Daarom stellen we
gekoppelde gegevens voor in een \textbf{spreidingsdiagram}.
Elk punt stelt één waarneming voor; alle punten samen vormen een
puntenwolk.

In dit hoofdstuk leren we zo'n puntenwolk zorgvuldig lezen.
We onderzoeken de \textbf{richting}, de \textbf{vorm} en de
\textbf{sterkte} van een mogelijk verband en letten op
\textbf{uitschieters} en \textbf{clusters}.

We rekenen het verband nog niet samen tot één getal en stellen nog
geen regressiemodel op. Eerst leren we kijken, beschrijven en
argumenteren. Dat vormt de noodzakelijke basis voor correlatie en
lineaire regressie in de volgende hoofdstukken.
\end{abstract}

\begin{document}

% FVD_CHAPTER_DATA_START
% Automatisch gegenereerd uit index.tex.
% Niet handmatig aanpassen.
\fvdchapterdata{3}{Verbanden onderzoeken met data}{1}
% FVD_CHAPTER_DATA_END

\maketitle

\begin{voorkennis}
  \item Je kan gegevens aflezen uit een tabel en een grafiek.
  \item Je kent het cartesisch assenstelsel en kan een punt $(x,y)$ plaatsen.
  \item Je kan grootheden, waarden en eenheden van elkaar onderscheiden.
  \item Je kan stijgend en dalend gedrag in een grafiek herkennen.
  \item Je weet dat meetgegevens een beschrijving van de werkelijkheid vormen.
\end{voorkennis}

\begin{leerdoelen}
  \item Je kan twee numerieke grootheden in een dataset herkennen en correct benoemen.
  \item Je kan uitleggen wat één punt in een spreidingsdiagram in de context betekent.
  \item Je kan gekoppelde gegevens voorstellen in een spreidingsdiagram.
  \item Je kan een puntenwolk beschrijven volgens richting, vorm en sterkte.
  \item Je kan een mogelijk positief of negatief verband herkennen.
  \item Je kan lineair en niet-lineair gedrag visueel van elkaar onderscheiden.
  \item Je kan uitschieters en clusters herkennen en uitleggen waarom ze aandacht verdienen.
  \item Je kan op basis van een spreidingsdiagram een voorzichtige, contextgebonden conclusie formuleren.
  \item Je kan uitleggen waarom een spreidingsdiagram op zichzelf geen causaliteit bewijst.
\end{leerdoelen}


% ============================================================
\FVDsection{Waarom twee grootheden samen onderzoeken?}
% ============================================================

In wetenschappen meten we vaak niet één, maar meerdere grootheden
bij hetzelfde object, dezelfde persoon of hetzelfde experiment.

Denk bijvoorbeeld aan:

\begin{itemize}
  \item de massa en de lengte van een dier;
  \item de hoogte en de gemiddelde temperatuur van een meetstation;
  \item de spanning en de stroomsterkte in een elektrische component;
  \item de concentratie van een oplossing en de gemeten absorbantie;
  \item de buitentemperatuur en het energieverbruik van een gebouw.
\end{itemize}

Wanneer twee numerieke grootheden bij dezelfde waarnemingen worden
gemeten, kunnen we onderzoeken of hun waarden \textbf{samen veranderen}.

\begin{opmerking}
Een verband onderzoeken begint niet met een formule.

We beginnen met drie vragen:

\[
\important{
\text{Wat werd gemeten?}
\quad
\text{Hoe zijn de gegevens gekoppeld?}
\quad
\text{Welk patroon zien we?}
}
\]
\end{opmerking}


\FVDsubsection{Gekoppelde gegevens}

Beschouw een eenvoudig experiment waarbij van acht metalen staven
de lengte en de massa worden gemeten.

\[
\begin{array}{c|cccccccc}
\text{staaf} & 1&2&3&4&5&6&7&8\\ \hline
\text{lengte }x\;(\mathrm{cm})
&10&14&18&22&26&30&34&38\\
\text{massa }y\;(\mathrm{g})
&24&32&43&50&61&69&81&88
\end{array}
\]

De waarden in dezelfde kolom horen bij \textbf{dezelfde staaf}.
We mogen de eerste waarneming dus schrijven als

\[
(10,24),
\]

de tweede als

\[
(14,32),
\]

enzovoort.

\begin{definitie}
Een dataset met twee numerieke grootheden bestaat uit
\textbf{gekoppelde waarnemingen}

\[
(x_1,y_1),(x_2,y_2),\ldots,(x_n,y_n).
\]

Bij elke waarneming horen dus twee waarden die samen gemeten zijn.
\end{definitie}

\begin{oefening}[Wat hoort bij elkaar?]{\oefB \ESS}
In een onderzoek worden voor verschillende dagen de buitentemperatuur
en het dagelijkse elektriciteitsverbruik van een koelinstallatie gemeten.

\begin{question}
Wat stelt één gegevenspaar $(x_i,y_i)$ voor?

\begin{multipleChoice}
  \choice{De temperatuur van één dag en het verbruik van een andere dag.}
  \choice[correct]{De temperatuur en het verbruik die bij dezelfde dag horen.}
  \choice{Twee opeenvolgende temperaturen.}
  \choice{Twee opeenvolgende verbruiken.}
\end{multipleChoice}
\end{question}

\begin{question}
Waarom mogen we de waarden niet willekeurig door elkaar mengen?

\begin{oplossing}
Omdat de koppeling tussen beide metingen essentieel is.
We willen onderzoeken hoe de twee grootheden bij dezelfde waarneming
samen voorkomen.
\end{oplossing}
\end{question}
\end{oefening}


% ============================================================
\FVDsection{Van tabel naar spreidingsdiagram}
% ============================================================

Een tabel bewaart de meetwaarden nauwkeurig, maar een patroon is er
niet altijd gemakkelijk in te herkennen.

Daarom zetten we ieder gegevenspaar $(x_i,y_i)$ als een punt in een
assenstelsel.

\begin{definitie}
Een \textbf{spreidingsdiagram} of \textbf{puntenwolk} is een grafische
voorstelling van gekoppelde numerieke gegevens.

Elk gegevenspaar

\[
(x_i,y_i)
\]

wordt voorgesteld door één punt in een assenstelsel.
\end{definitie}

Voor de metalen staven krijgen we dus acht punten. De punten hoeven
niet exact op een rechte of kromme te liggen. Het gaat om
\textbf{meetgegevens}; natuurlijke variatie en meetonzekerheid zijn
normaal.


\FVDsubsection{Wat stelt één punt voor?}

Een punt in een spreidingsdiagram heeft altijd een betekenis in de
context.

Het punt

\[
(26,61)
\]

betekent in het voorbeeld:

\begin{quote}
Een staaf met een lengte van $26\,\mathrm{cm}$ heeft een gemeten
massa van $61\,\mathrm{g}$.
\end{quote}

\begin{opmerking}
Lees een punt nooit alleen als ``$(26,61)$''.

Een goede interpretatie bevat:

\begin{itemize}
  \item de twee grootheden;
  \item de twee waarden;
  \item de eenheden;
  \item de betekenis in de context.
\end{itemize}
\end{opmerking}


\FVDsubsection{Welke grootheid komt op welke as?}

Wanneer de context een natuurlijke richting suggereert, plaatsen we
de grootheid die als mogelijke verklarende of instelbare grootheid
wordt gebruikt meestal op de horizontale as.

De andere grootheid komt dan op de verticale as.

Bijvoorbeeld:

\[
\text{hoogte}
\longrightarrow
\text{temperatuur}
\]

of

\[
\text{concentratie}
\longrightarrow
\text{absorbantie}.
\]

\begin{opmerking}
De keuze van de assen \textbf{bewijst geen oorzaak-gevolgrelatie}.
Ze helpt ons alleen om de gegevens op een zinvolle manier te organiseren.
Of een verband werkelijk causaal is, vraagt bijkomende argumenten en
onderzoek.
\end{opmerking}


\begin{oefening}[De assen kiezen]{\oefB \ESS}
Een leerling onderzoekt bij een veer hoe de uitrekking verandert
wanneer verschillende massa's worden opgehangen.

\begin{question}
Welke grootheid plaats je hier het meest natuurlijk op de horizontale as?

\begin{multipleChoice}
  \choice[correct]{de opgehangen massa}
  \choice{de uitrekking}
\end{multipleChoice}
\end{question}

\begin{question}
Welke informatie moet bij de assen vermeld worden?

\begin{oplossing}
De naam of het symbool van de grootheid en de gebruikte eenheid.
\end{oplossing}
\end{question}
\end{oefening}


% ============================================================
\FVDsection{Een puntenwolk systematisch lezen}
% ============================================================

Een puntenwolk beoordelen we niet op basis van één punt.
We bekijken het \textbf{globale patroon}.

Daarbij gebruiken we vier aandachtspunten:

\[
\important{
\text{richting}
\quad
\text{vorm}
\quad
\text{sterkte}
\quad
\text{bijzonderheden}
}
\]

Die vier woorden vormen vanaf nu onze vaste leesstrategie.


% ------------------------------------------------------------
\FVDsubsection{1. Richting}
% ------------------------------------------------------------

\begin{definitie}
We spreken van een \textbf{positief verband} wanneer grotere waarden
van de ene grootheid in het algemeen samengaan met grotere waarden
van de andere grootheid.

We spreken van een \textbf{negatief verband} wanneer grotere waarden
van de ene grootheid in het algemeen samengaan met kleinere waarden
van de andere grootheid.
\end{definitie}

Positief en negatief zeggen dus iets over de \textbf{richting} van
de puntenwolk, niet over hoe sterk het verband is.

\begin{voorbeeld}[Twee wetenschappelijke situaties]
Bij eenzelfde soort metalen staaf verwachten we doorgaans een positief
verband tussen lengte en massa.

Bij meetstations in een berggebied verwachten we doorgaans een negatief
verband tussen hoogte en temperatuur: hoger gelegen stations zijn
gemiddeld kouder.
\end{voorbeeld}


% ------------------------------------------------------------
\FVDsubsection{2. Vorm}
% ------------------------------------------------------------

Niet ieder verband is lineair.

\begin{definitie}
Een puntenwolk vertoont een \textbf{ongeveer lineair verband} wanneer
de punten globaal rond een rechte band liggen.

Bij een \textbf{niet-lineair verband} volgt de puntenwolk duidelijk
een andere vorm, bijvoorbeeld een kromme.
\end{definitie}

Dat onderscheid is belangrijk. In het volgende hoofdstuk zullen we
een getal invoeren dat specifiek de sterkte van een
\textbf{lineair} verband beschrijft.

\begin{opmerking}
Een puntenwolk zonder duidelijke rechte richting kan toch een sterk
verband bevatten.

Bijvoorbeeld: wanneer $y$ ongeveer afhangt van $x^2$, kan de puntenwolk
een duidelijke U-vorm hebben. Het verband is dan sterk, maar niet lineair.
\end{opmerking}


% ------------------------------------------------------------
\FVDsubsection{3. Sterkte}
% ------------------------------------------------------------

Wanneer de punten dicht rond een herkenbaar patroon liggen, spreken
we informeel van een sterker verband.

Wanneer de punten veel verder verspreid liggen, is het verband zwakker.

\begin{opmerking}
In dit hoofdstuk beoordelen we de sterkte nog \textbf{visueel}.

In het volgende hoofdstuk leren we de sterkte en richting van een
lineair verband samenvatten met de correlatiecoëfficiënt.
\end{opmerking}


% ------------------------------------------------------------
\FVDsubsection{4. Bijzonderheden: uitschieters en clusters}
% ------------------------------------------------------------

Een globale trend is niet het hele verhaal.

\begin{definitie}
Een \textbf{uitschieter} is een waarneming die opvallend ver van het
algemene patroon van de andere gegevens ligt.
\end{definitie}

Een uitschieter kan verschillende oorzaken hebben:

\begin{itemize}
  \item een meet- of invoerfout;
  \item een uitzonderlijke maar correcte waarneming;
  \item een andere experimentele situatie;
  \item een aanwijzing dat ons model te eenvoudig is.
\end{itemize}

We verwijderen een uitschieter daarom nooit automatisch.

\begin{definitie}
Een \textbf{cluster} is een groep waarnemingen die in een bepaald
gebied van het spreidingsdiagram duidelijk samen ligt.
\end{definitie}

Clusters kunnen erop wijzen dat verschillende groepen in de dataset
door elkaar voorkomen.

Denk bijvoorbeeld aan metingen bij twee verschillende materialen,
twee leeftijdsgroepen of twee experimentele omstandigheden.


% ============================================================
\FVDsection{Van kijken naar beschrijven}
% ============================================================

Een goede analyse bestaat niet uit één woord zoals ``positief'' of
``sterk''.

Gebruik bij een nieuwe puntenwolk bij voorkeur deze volgorde:

\begin{enumerate}
  \item \textbf{Context:} welke twee grootheden worden vergeleken?
  \item \textbf{Richting:} positief, negatief of geen duidelijke richting?
  \item \textbf{Vorm:} ongeveer lineair of duidelijk niet-lineair?
  \item \textbf{Sterkte:} liggen de punten dicht of ruim rond het patroon?
  \item \textbf{Bijzonderheden:} zijn er uitschieters, clusters of andere opvallende kenmerken?
  \item \textbf{Conclusie:} wat suggereert de puntenwolk in deze context?
\end{enumerate}

\begin{voorbeeld}[Hoogte en temperatuur]
Stel dat meetstations op verschillende hoogtes worden vergeleken en
de puntenwolk duidelijk van linksboven naar rechtsonder loopt.

Een goede beschrijving is bijvoorbeeld:

\begin{quote}
De gegevens suggereren een vrij sterk negatief en ongeveer lineair
verband tussen hoogte en temperatuur. Hoger gelegen meetstations
hebben in deze dataset doorgaans een lagere temperatuur. Eén station
wijkt opvallend af van het globale patroon.
\end{quote}

Merk op wat we \textbf{niet} schrijven:

\begin{quote}
Hoogte veroorzaakt hier zeker de gemeten temperatuur.
\end{quote}

Een spreidingsdiagram toont samenhang in de gegevens, maar bewijst
op zichzelf geen causaliteit.
\end{voorbeeld}


\begin{oefening}[Van losse woorden naar een analyse]{\oefU \ESS}
Een leerling bekijkt een spreidingsdiagram van buitentemperatuur en
dagelijks gasverbruik van een schoolgebouw. De punten liggen vrij dicht
rond een dalende rechte band, behalve één punt dat veel hoger ligt.

\begin{question}
Welke richting heeft het verband?

\begin{multipleChoice}
  \choice{positief}
  \choice[correct]{negatief}
  \choice{geen duidelijke richting}
\end{multipleChoice}
\end{question}

\begin{question}
Welke vorm lijkt geschikt als eerste beschrijving?

\begin{multipleChoice}
  \choice[correct]{ongeveer lineair}
  \choice{duidelijk U-vormig}
  \choice{volledig willekeurig}
\end{multipleChoice}
\end{question}

\begin{question}
Wat moet je met het afwijkende punt doen?

\begin{multipleChoice}
  \choice{Het onmiddellijk verwijderen.}
  \choice{Het negeren omdat de andere punten wel goed liggen.}
  \choice[correct]{Het onderzoeken en nagaan waarom die waarneming afwijkt.}
\end{multipleChoice}
\end{question}

\begin{question}
Formuleer een volledige conclusie in woorden.

\begin{oplossing}
De gegevens suggereren een vrij sterk negatief en ongeveer lineair
verband tussen buitentemperatuur en gasverbruik: op warmere dagen is
het gasverbruik doorgaans lager. Eén waarneming wijkt duidelijk af en
verdient verder onderzoek. Uit deze puntenwolk alleen volgt nog geen
bewijs voor causaliteit.
\end{oplossing}
\end{question}
\end{oefening}


% ============================================================
\FVDsection{Een spreidingsdiagram maken met ICT}
% ============================================================

Bij grotere datasets gebruiken we ICT om de puntenwolk te tekenen.
Dat bespaart tijd, maar de software neemt het denkwerk niet over.

Een correcte werkwijze is:

\begin{enumerate}
  \item controleer welke twee kolommen gekoppelde gegevens bevatten;
  \item kies welke grootheid op de horizontale en verticale as komt;
  \item maak een spreidingsdiagram;
  \item geef beide assen een betekenisvolle naam en eenheid;
  \item controleer de schaal;
  \item beschrijf daarna richting, vorm, sterkte en bijzonderheden.
\end{enumerate}

\begin{opmerking}
Een lijngrafiek en een spreidingsdiagram zijn niet hetzelfde.

Bij een spreidingsdiagram worden afzonderlijke gekoppelde waarnemingen
als punten weergegeven. De punten hoeven niet in meetvolgorde met
lijnstukken verbonden te worden.
\end{opmerking}


% ============================================================
\FVDsection{Grafieken kritisch bekijken}
% ============================================================

Ook een correct getekende puntenwolk kan verkeerd geïnterpreteerd worden.

Let daarom op:

\begin{itemize}
  \item de schaal van beide assen;
  \item het bereik van de gemeten waarden;
  \item het aantal waarnemingen;
  \item mogelijke uitschieters;
  \item mogelijke clusters;
  \item de vraag of het globale patroon werkelijk lineair lijkt.
\end{itemize}

\begin{voorbeeld}[Een te snelle conclusie]
Een onderzoek bevat slechts vijf meetpunten. Vier punten liggen dicht
bij elkaar en één punt ligt ver weg. Door dat ene punt lijkt een sterk
stijgend verband te ontstaan.

Dan is ``er is een sterk positief verband'' een te snelle conclusie.
De uitschieter heeft mogelijk een zeer grote invloed op het beeld.
We moeten eerst onderzoeken wat er met die waarneming aan de hand is.
\end{voorbeeld}


% ============================================================
\FVDsection{Samengevat}
% ============================================================

Wanneer twee numerieke grootheden bij dezelfde waarnemingen worden
gemeten, kunnen we hun samenhang onderzoeken met een spreidingsdiagram.

\[
\important{
(x_i,y_i)
\longrightarrow
\text{punt}
\longrightarrow
\text{puntenwolk}
\longrightarrow
\text{patroon}
\longrightarrow
\text{voorzichtige conclusie}
}
\]

Bij het lezen van een puntenwolk onderzoeken we steeds:

\[
\boxed{
\text{richting}
\quad
\text{vorm}
\quad
\text{sterkte}
\quad
\text{bijzonderheden}
}
\]

Een positief verband betekent dat grotere waarden doorgaans samengaan
met grotere waarden. Bij een negatief verband gaan grotere waarden
doorgaans samen met kleinere waarden.

Een verband kan ongeveer lineair of duidelijk niet-lineair zijn.
Uitschieters en clusters mogen niet zomaar genegeerd worden.

Ten slotte geldt:

\[
\important{
\text{een spreidingsdiagram toont een patroon in gegevens,
maar bewijst geen causaliteit}
}
\]

In het volgende hoofdstuk gaan we een ongeveer lineair verband
nauwkeuriger beschrijven met de \textbf{correlatiecoëfficiënt}.

\end{document}
